\begin{helper}
Let \(k,u_R,u_B,N_R,N_B\) be natural numbers and let \(p,\varepsilon_1,a\) be
real numbers.  Assume
\[
0\le\varepsilon_1\le1,\qquad 0\le p\le\frac12,\qquad
u_B\le u_R,
\]
and
\[
N_R\le 3\varepsilon_1k^2,\qquad N_B\le k^2.
\]
Then
\[
\binom{N_R}{u_R}p^{u_R}
\binom{N_B}{u_B}p^{u_B}
(1-p)^{\max(0,a-u_R-u_B)}
\le
\exp\!\left(8\sqrt{\varepsilon_1}\,pk^2-pa\right).
\]
\end{helper}
\begin{proof}
First handle \(p=0\).  If \(u_R+u_B>0\), one of the factors \(p^{u_R}\) or
\(p^{u_B}\) is zero, so the left-hand side is \(0\).  If
\(u_R+u_B=0\), then \(u_R=u_B=0\), both binomial coefficients are \(1\), and
the absence factor is at most \(1\): when \(a\ge0\) it is
\((1-0)^a=1\), and when \(a<0\) the displayed exponent is \(0\).  The
right-hand side is \(\exp(0)=1\).  Thus the estimate holds when \(p=0\).
Assume from now on that \(p>0\), and only now set
\[
q=2p.
\]
Then \(0<q\le1\).

Next split the edge case \(\varepsilon_1=0\) before taking any logarithm
involving \(3e\varepsilon_1\).  The inequality
\(N_R\le3\varepsilon_1k^2=0\) forces \(N_R=0\).  If \(u_R>0\), then
\(\binom{N_R}{u_R}=0\), so the left-hand side is \(0\).  If \(u_R=0\), then
\(u_B=0\) because \(u_B\le u_R\), and the support part is \(1\).  The desired
inequality reduces to
\[
(1-p)^{\max(0,a)}\le \exp(-pa).
\]
If \(a\ge0\), this follows from \(1-p\le\exp(-p)\).  If \(a<0\), the left
side is \(1\), while \(\exp(-pa)\ge1\) because \(p>0\).  Hence the theorem is
proved when \(\varepsilon_1=0\).  We may therefore assume
\[
0<\varepsilon_1\le1. \tag{1}
\]

We prove the two-support estimate with \(q=2p\).  Put \(K=k^2\).  If
\(K=0\), then \(N_R=N_B=0\).  The support product is \(0\) unless
\(u_R=u_B=0\), and the case \(u_R=u_B=0\) is handled by the same absence
argument used above.  Thus assume \(K>0\).

If \(u_R=0\), then \(u_B=0\), and the support estimate below is simply
\[
1\le \exp(8\sqrt{\varepsilon_1}pK),
\]
which follows from \(p\ge0\) and \(\varepsilon_1\ge0\).  Now suppose
\(u_R>0\).  Define
\[
s=\frac{u_R}{qK},\qquad t=\frac{u_B}{qK}.
\]
Then \(s>0\), \(0\le t\le s\), and the convention is that
\(t\log(e/t)=0\) when \(t=0\).  Using
\(\binom Nx\le(eN/x)^x\) for \(x>0\), the hypotheses
\[
N_R\le3\varepsilon_1K,\qquad N_B\le K,
\]
and the same estimate for the \(u_B\)-term when \(u_B>0\), we get
\[
\binom{N_R}{u_R}q^{u_R}
\binom{N_B}{u_B}q^{u_B}
\le
\exp\!\left(qK\left[
s\log\!\left(\frac{3e\varepsilon_1}{s}\right)
+t\log\!\left(\frac e t\right)\right]\right). \tag{2}
\]
No logarithm at \(t=0\) is used; the \(u_B=0\) binomial factor is \(1\), which
is represented by the zero value of \(t\log(e/t)\).

We now prove the elementary calculus bound used in (2).  Let
\(0<\delta\le1\), \(s>0\), and \(0\le t\le s\).  Write \(r=t/s\).  Then
\[
s\log\!\left(\frac{3e\delta}{s}\right)
+t\log\!\left(\frac e t\right)
=
s\left[
\log\!\left(\frac{3e\delta}{s}\right)
+r\log\!\left(\frac e{rs}\right)\right],        \tag{3}
\]
again with the limiting value \(r\log(e/(rs))=0\) at \(r=0\).  For fixed
\(s\), the derivative in \(r\) of the bracketed expression is
\(\log(1/(rs))\).  Hence, on \(0\le r\le1\), the maximum occurs at
\(r=1\) if \(s\le1\), and at \(r=1/s\) if \(s\ge1\).

If \(s\le1\), (3) is at most
\[
s\log\!\left(\frac{3e^2\delta}{s^2}\right).
\]
The derivative of this function of \(s\) is
\(\log(3\delta/s^2)\), so its maximum on \(0<s\le1\) is no larger than the
unrestricted maximum at \(s=\sqrt{3\delta}\), namely
\[
2\sqrt{3\delta}\le4\sqrt\delta. \tag{4}
\]
If \(s\ge1\), (3) is at most
\[
s\log\!\left(\frac{3e\delta}{s}\right)+1. \tag{5}
\]
When \(\delta\le1/3\), the derivative of (5) is
\(\log(3\delta/s)\le0\) for \(s\ge1\), so the maximum is attained at
\(s=1\) and equals \(1+\log(3e\delta)\).  The function
\[
4\sqrt\delta-1-\log(3e\delta)
\]
is nonnegative on \(0<\delta\le1/3\) (its minimum in this interval occurs at
\(\delta=1/4\), where it is \(1-\log(3e/4)>0\), or at an endpoint).  Thus
(5) is at most \(4\sqrt\delta\) in this case.  When
\(\delta\ge1/3\), the maximum of (5) on \(s\ge1\) occurs at \(s=3\delta\) and
has value \(1+3\delta\).  For \(1/3\le\delta\le1\),
\[
1+3\delta\le4\sqrt\delta
\]
because \(3x^2-4x+1=(3x-1)(x-1)\le0\) for
\(x=\sqrt\delta\in[1/\sqrt3,1]\).  Combining this with (4) proves
\[
s\log\!\left(\frac{3e\delta}{s}\right)
+t\log\!\left(\frac e t\right)
\le4\sqrt\delta. \tag{6}
\]
Applying (6) with \(\delta=\varepsilon_1\) in (2), and using \(q=2p\), gives
\[
\binom{N_R}{u_R}q^{u_R}
\binom{N_B}{u_B}q^{u_B}
\le
\exp(4qK\sqrt{\varepsilon_1})
=
\exp(8\sqrt{\varepsilon_1}\,p k^2). \tag{7}
\]

It remains to replace the \(q\)-support factors by the original \(p\)-support
factors and the absence term.  Let
\[
U=u_R+u_B,\qquad c=a-u_R-u_B.
\]
If \(c\ge0\), then the absence factor is \((1-p)^c\).  Since
\(1-p\le\exp(-p)\),
\[
(1-p)^c\le\exp(-pc).
\]
Also \(p\le1/2\) implies \(\exp(p)\le2\), so
\[
p^U\exp(pU)\le(2p)^U=q^U.
\]
Multiplying the last two displays gives
\[
p^U(1-p)^c\le q^U\exp(-pa). \tag{8}
\]

If \(c<0\), the absence exponent in the statement is \(0\).  We claim
\[
p^U\le q^U\exp(-pa). \tag{9}
\]
If \(a\le0\), then \(\exp(-pa)\ge1\) and (9) follows from \(p\le q\).  If
\(0<a<U\), then \(p\le1/2<\log2\), so
\[
2^{-U}\le\exp(-pU)<\exp(-pa).
\]
Since \(p^U=q^U2^{-U}\), (9) follows.

In both cases, the product of the original support powers and the displayed
absence factor is bounded by \(q^{u_R+u_B}\exp(-pa)\).  Multiplying this with
the binomial part and using (7) yields
\[
\binom{N_R}{u_R}p^{u_R}
\binom{N_B}{u_B}p^{u_B}
(1-p)^{\max(0,a-u_R-u_B)}
\le
\exp(8\sqrt{\varepsilon_1}pk^2)\exp(-pa).
\]
Adding the exponents gives the required
\[
\exp\!\left(8\sqrt{\varepsilon_1}pk^2-pa\right).
\]
\end{proof}
